% Source: https://artofproblemsolving.com/community/c3h1644585_system_of_equations
% Stable identity: aops-topic-1644585-q01
The first equation says $27xy=71$, so $xy=\dfrac{71}{27}$. From the second, $x=1+\dfrac{3y}{2}$. Substituting,
\[
\left(1+\dfrac{3y}{2}\right)y=\dfrac{71}{27}\iff\dfrac{3}{2}y^2+y-\dfrac{71}{27}=0\iff81y^2+54y-142=0 .
\]
The discriminant is $54^2+4\cdot81\cdot142=2916+46008=48924=18^2\cdot151$, so
\[
y=\dfrac{-54\pm18\sqrt{151}}{162}=\dfrac{-3\pm\sqrt{151}}{9},\qquad x=1+\dfrac{3y}{2}=\dfrac{3\pm\sqrt{151}}{6}
\]
(with the same sign). The two solutions are
\[
\boxed{(x,y)=\left(\dfrac{3+\sqrt{151}}{6},\ \dfrac{-3+\sqrt{151}}{9}\right)\ \text{and}\ \left(\dfrac{3-\sqrt{151}}{6},\ \dfrac{-3-\sqrt{151}}{9}\right)}.
\]
Check for the first: $xy=\dfrac{(\sqrt{151}+3)(\sqrt{151}-3)}{54}=\dfrac{142}{54}=\dfrac{71}{27}$, and $2x-3y=\dfrac{3+\sqrt{151}}{3}-\dfrac{-3+\sqrt{151}}{3}=2$.
